\documentclass[nonacm]{acmart}

\AtBeginDocument{%
  \providecommand\BibTeX{{%
    Bib\TeX}}}

\setcopyright{none}

\usepackage{array}
\usepackage{mathtools}
\usepackage{stmaryrd}
\usepackage{mathpartir}
\usepackage{listings}
\usepackage{xcolor}

\lstset{
  basicstyle=\ttfamily\small,
  breaklines=true,
  columns=fullflexible,
  keepspaces=true,
  frame=single,
  showstringspaces=false
}

\newcommand{\letterm}[3]{\mathbf{let}\ \{#1\}=#2\ \mathbf{in}\ #3}
\newcommand{\as}{\mathrel{\mathbf{as}}}

\begin{document}

\title{Encodings of Existentials in System F}

\author{Runbang Yan}
\email{rbyan25@stu.pku.edu.cn}
\affiliation{%
  \institution{School of Computer Science, Peking University}
  \city{Beijing}
  \country{China}
}

\renewcommand{\shortauthors}{Yan}

\maketitle

\section{Evaluation of Encoded Existential Elimination}

\textbf{Claim.}
Under the standard System F encoding of existential types,
\[
\{\exists X, T\}
\triangleq
\forall S.\,(\forall X.\,T \to S) \to S,
\]
we have
\[
\letterm{X,x}{(\{*T_{11}, v_{12}\} \as T_1)}{t_2}
\longrightarrow^{*}
[X \mapsto T_{11}][x \mapsto v_{12}]t_2.
\]

\textbf{Proof.}
Assume that $T_1 = \{\exists X,T_{12}\}$, and suppose that $t_2$ has type $S$ in the context extended with the abstract type variable $X$ and the term variable $x:T_{12}$.
We also assume the usual side condition for existential elimination, namely that $X$ does not occur free in the result type $S$.

To avoid confusion between the witness type and the answer type of the encoding, we alpha-rename the bound answer type variable in the package encoding to $R$.
Thus
\[
\{*T_{11}, v_{12}\} \as T_1
\triangleq
\lambda R.\,\lambda f:(\forall X.\,T_{12}\to R).\,f\,[T_{11}]\,v_{12}.
\]
By the encoding of \textsf{let} for existentials,
\[
\letterm{X,x}{t_1}{t_2}
\triangleq
t_1\,[S]\,(\lambda X.\,\lambda x:T_{12}.\,t_2).
\]
Therefore,
\[
\begin{aligned}
&\letterm{X,x}{(\{*T_{11}, v_{12}\} \as T_1)}{t_2}
\\
&\quad =
(\lambda R.\,\lambda f:(\forall X.\,T_{12}\to R).\,f\,[T_{11}]\,v_{12})
\,[S]\,
(\lambda X.\,\lambda x:T_{12}.\,t_2)
\\
&\quad \longrightarrow
(\lambda f:(\forall X.\,T_{12}\to S).\,f\,[T_{11}]\,v_{12})
\,(\lambda X.\,\lambda x:T_{12}.\,t_2)
\\
&\quad \longrightarrow
(\lambda X.\,\lambda x:T_{12}.\,t_2)\,[T_{11}]\,v_{12}
\\
&\quad \longrightarrow
(\lambda x:[X\mapsto T_{11}]T_{12}.\,[X\mapsto T_{11}]t_2)\,v_{12}
\\
&\quad \longrightarrow
[x\mapsto v_{12}]([X\mapsto T_{11}]t_2)
\\
&\quad =
[X\mapsto T_{11}][x\mapsto v_{12}]t_2.
\end{aligned}
\]
The first and third reduction steps are type-level beta-reductions, and the second and fourth reduction steps are ordinary term-level beta-reductions.
All substitutions are capture-avoiding.
Hence
\[
\letterm{X,x}{(\{*T_{11}, v_{12}\} \as T_1)}{t_2}
\longrightarrow^{*}
[X \mapsto T_{11}][x \mapsto v_{12}]t_2.
\]
\[
\Box
\]



\end{document}
